% Bản nháp chương 5 theo dàn bài mới, viết với bộ mẫu ở style/mau/.
% Chưa nối vào report.tex.
% ===========================================================================
\chapter{Cài đặt}
\label{sec:algorithms}

Năm thuật toán so sánh ở chương~\ref{sec:compare} đều do nhóm tự cài, dùng chung
một chữ ký hàm và trả về chung một cấu trúc lịch sử. Chương này mô tả cập nhật
của từng thuật toán cùng chi phí mỗi vòng lặp của nó, và trước hết là phần đo,
vì mọi con số giây trong báo cáo chỉ so được với nhau khi phép đo giống nhau.

\section{Giao diện chung và cách đo}
\label{sec:recorder}

Mọi thuật toán nhận cùng một bộ tham số gồm bài toán, điểm khởi đầu, ngân sách
vòng lặp, quy tắc chọn độ dài bước, sai số dừng, tần suất ghi lịch sử và hạt
giống ngẫu nhiên, rồi trả về một đối tượng chứa dãy vòng lặp, dãy giá trị hàm
mục tiêu, dãy chuẩn gradient, dãy thời gian tích lũy, dãy số lần truy cập dữ
liệu và trạng thái dừng. Nhờ vậy phần vẽ hình dùng lại được cho mọi phương pháp,
và một cấu hình mới không kéo theo một nhánh riêng trong mã vẽ.

Phần đo thời gian có ba chi tiết quyết định tính đúng của mọi biểu đồ theo trục
giây. Thứ nhất, đồng hồ dừng trước khi tính $f(w_k)$ và chuẩn gradient rồi chạy
lại sau đó, vì hai phép tính ấy tốn ngang một vòng lặp gradient nên ghi lịch sử
dày sẽ thổi phồng thời gian nếu tính vào. Thứ hai, bộ đếm số lần đánh giá hàm và
gradient được khôi phục sau mỗi lần ghi, nếu không thì cột chi phí ở
chương~\ref{sec:mechanisms} sẽ đếm cả phần theo dõi. Thứ ba, mỗi lần đo chạy vài
vòng khởi động rồi bỏ kết quả đi, vì lần gọi thư viện đại số tuyến tính đầu tiên
mang theo chi phí khởi tạo. Cấu hình máy, phiên bản thư viện và seed dùng chung
ghi ở phụ lục~\ref{app:config}.

Ba chi tiết ấy loại được chi phí của phép theo dõi nhưng không loại được ảnh
hưởng của nó lên phần cứng. Phép đo RMSE ở chương~\ref{sec:enough} quét ma trận
kiểm tra \SI{85}{\mebi\byte} xen giữa các vòng lặp và đẩy ma trận huấn luyện
\SI{342}{\mebi\byte} ra khỏi cache, làm lần chạy chậm đi 13\% dù đồng hồ đã dừng
đúng lúc. Vì thế các thí nghiệm so sánh thời gian chạy không bật phép theo dõi
ấy, và chương~\ref{sec:enough} chỉ báo số vòng lặp.

Cuối cùng, một lần chạy dừng khi chuẩn gradient tương đối xuống dưới ngưỡng, khi
hết ngân sách vòng lặp, khi hàm mục tiêu bùng lên quá $10^{12}$ lần giá trị đầu,
hoặc khi nó không cải thiện quá một sai số tương đối $10^{-13}$ trong mười lần
ghi liên tiếp. Tiêu chí cuối cùng có mặt vì một lý do đo được: trước khi thêm nó,
một lần chạy backtracking đạt sai số $10^{-15}$ ở vòng lặp 115 rồi chạy tiếp tới
vòng 250, và 98,7\% trong 72 giây của lần chạy ấy là thời gian chết. Tiêu chí này
không áp cho các phương pháp ngẫu nhiên, vì hàm mục tiêu của chúng không đơn điệu
và chỗ nằm ngang chính là kết quả cần trình bày.

\section{Gradient descent}
\label{sec:impl-gd}

\begin{algorithm}[tp]
  \caption{Gradient descent}
  \label{alg:gd}
  \begin{algorithmic}[1]
    \Require objective $f$, initial point $w_0$, step size rule
    \For{$k = 0, 1, 2, \dots$}
      \State $g_k \gets \nabla f(w_k)$
      \State choose $t_k$ (fixed, or by Armijo backtracking)
      \State $w_{k+1} \gets w_k - t_k g_k$
    \EndFor
  \end{algorithmic}
\end{algorithm}

Thuật toán~\ref{alg:gd} dùng đúng một thông tin về hàm mục tiêu là gradient tại
điểm hiện tại, nên mỗi vòng lặp tốn $\bigO(nd)$, bằng một lần quét toàn bộ ma
trận thiết kế. Với hàm bậc hai lồi mạnh, phương pháp hội tụ khi $0 < t < 2/L$ và
độ dài bước tối ưu là $t = 2/(L+\mu)$~\cite{nesterov2018lectures}. Hai mốc ấy đến
thẳng từ \eqref{eq:lipschitz}: sai số theo hướng riêng ứng với trị riêng
$\lambda_i$ nhân với $\abs{1 - t\lambda_i}$ mỗi vòng, nên điều kiện hội tụ là hệ
số ấy nhỏ hơn 1 với mọi $\lambda_i \in [\mu, L]$, và độ dài bước tối ưu là giá
trị làm cân bằng hai đầu khoảng. Chương~\ref{sec:mechanisms} đo cả hai mốc.

\section{Backtracking theo điều kiện Armijo}
\label{sec:impl-armijo}

Xuất phát từ $t = t_0$, thu nhỏ $t \leftarrow \beta t$ cho tới
khi~\cite{nocedal2006numerical}
%
\begin{equation}
  \label{eq:armijo}
  f\bigl( w + t p \bigr) \le f(w) + \alpha \, t \, \langle \nabla f(w), p \rangle ,
\end{equation}
%
với $p$ là hướng giảm. Khi $p = -\nabla f(w)$, điều kiện trở thành
$f(w - t \nabla f(w)) \le f(w) - \alpha t \norm{\nabla f(w)}_{2}^{2}$.

Điều kiện \eqref{eq:armijo} không nhận mọi bước làm hàm giảm mà đòi mức giảm đạt
ít nhất tỉ lệ $\alpha$ của mức giảm mà xấp xỉ tuyến tính hứa hẹn. Đòi hỏi ấy cần
thiết vì một bước rất nhỏ luôn làm $f$ giảm chút ít, và nhận nó thì thuật toán bò
mãi mà không tới đâu. Hai tham số vì thế đóng hai vai khác nhau: $\alpha$ định
nghĩa thế nào là giảm đủ, còn $\beta$ định chi phí tìm kiếm vì nó quyết định mỗi
lần thử thu nhỏ bước bao nhiêu. Chương~\ref{sec:mechanisms} cho thấy chỉ $\beta$
đổi được số lần đánh giá hàm mỗi vòng lặp.

\section{Stochastic gradient descent}
\label{sec:impl-sgd}

\begin{equation}
  \label{eq:sgd}
  g_k = \frac{1}{\abs{\Batch_k}} \sum_{i \in \Batch_k}
        \bigl( x_i\trans w_k - y_i \bigr) x_i + \lambda w_k,
  \qquad
  w_{k+1} = w_k - \eta_k g_k .
\end{equation}

Ở đây $\Batch_k$ trong \eqref{eq:sgd} là một lô chỉ số lấy ngẫu nhiên, nên $g_k$
là ước lượng không chệch của gradient đầy đủ và mỗi vòng lặp chỉ tốn $\bigO(Bd)$ thay vì
$\bigO(nd)$. Đổi lại, $g_k$ mang phương sai không tắt theo $k$, và
chương~\ref{sec:mechanisms} cho thấy chính phương sai ấy quyết định SGD dừng ở
đâu.

Nhóm thử bốn quy tắc chọn độ dài bước~\cite{bottou2018optimization}
%
\begin{equation}
  \label{eq:schedules}
  \eta_k = \eta_0, \qquad
  \eta_k = \frac{\eta_0}{1 + \gamma k}, \qquad
  \eta_k = \frac{\eta_0}{\sqrt{k+1}}, \qquad
  \eta_k = \eta_0 \, 2^{-\lfloor k / p \rfloor} .
\end{equation}

Bốn quy tắc ở \eqref{eq:schedules} ấn định $\eta_k$ trước theo chỉ số vòng lặp và
không đọc gì từ giá trị hàm mục tiêu, nên chúng là quy tắc tất định chứ không
phải một dạng line search. Ba thuật toán còn lại của báo cáo đều có một biến thể
dùng backtracking, riêng SGD thì không, và lý do đo được. Điều kiện Armijo
\eqref{eq:armijo} đòi đánh giá $f$ trên toàn bộ tập huấn luyện ở mỗi lần thử độ
dài bước, trong khi với $B = 64$ một lượt duyệt dữ liệu gồm \num{2500} lần cập
nhật. Lấy \SI{23}{\milli\second} làm chi phí một lần quét toàn bộ dữ liệu thì
riêng phần line search cần khoảng \SI{58}{\second} cho một lượt duyệt, còn cả
lượt duyệt đo được chỉ mất \SI{0.087}{\second}, tức gấp hơn 600 lần. Con số
\SI{58}{\second} là ước lượng suy ra từ hai đại lượng đo được và chưa tính số lần
thu nhỏ bước trong mỗi lần gọi, nên chênh lệch thật chỉ có thể lớn hơn.

Trở ngại thứ hai nằm ở phần bảo đảm. Áp \eqref{eq:armijo} lên $g_k$ chỉ bảo đảm
giảm cho hàm mất mát của lô hiện tại chứ không cho $f$, nên một độ dài bước được
chấp nhận trên lô này vẫn có thể làm $f$ tăng. Vì vậy độ dài bước ban đầu lấy
theo hằng số trơn của lô, $\eta_0 = 1/L_B$.

\section{Accelerated gradient descent}
\label{sec:impl-agd}

\begin{equation}
  \label{eq:agd}
  y_k = w_k + \beta_k (w_k - w_{k-1}), \qquad
  w_{k+1} = y_k - t \nabla f(y_k) ,
\end{equation}
%
với $\beta = (\sqrt{\kappa} - 1)/(\sqrt{\kappa} + 1)$ khi biết
$\mu$~\cite{nesterov2018lectures}, hoặc $\beta_k = (k-1)/(k+2)$ khi
không~\cite{beck2009fista}.

So với thuật toán~\ref{alg:gd}, cập nhật \eqref{eq:agd} chỉ thêm một điểm nhớ là
$w_{k-1}$, nên chi phí mỗi vòng vẫn là $\bigO(nd)$ và phần thêm chỉ là một phép
cộng vector. Điểm khác nằm ở chỗ gradient được lấy tại điểm ngoại suy $y_k$ chứ
không tại $w_k$, tức thuật toán đọc hàm mục tiêu ở phía trước hướng đang đi. Cơ
chế khởi động lại đặt $k = 0$ mỗi khi $\nabla f(y_k)\trans (w_{k+1} - w_k) >
0$~\cite{odonoghue2015restart}, tức mỗi khi bước vừa đi ngược hướng giảm, và
chương~\ref{sec:mechanisms} đo tác dụng của nó.

\section{Phương pháp Newton}
\label{sec:impl-newton}

Bước Newton giải hệ $\nabla^{2} f(w_k) p_k = -\nabla f(w_k)$ bằng phân rã
Cholesky chứ không nghịch đảo ma trận~\cite{boyd2004convex}. Với hàm mục tiêu bậc
hai và $w_0 = 0$,
%
\begin{equation}
  \label{eq:newton-one-step}
  w_1 = w_0 - \bigl( \nabla^{2} f \bigr)^{-1} \nabla f(w_0)
      = \Bigl( \tfrac{1}{n} X\trans X + \lambda I \Bigr)^{-1}
        \Bigl( \tfrac{1}{n} X\trans y \Bigr) = \wstar ,
\end{equation}
%
nghĩa là phương pháp hội tụ sau đúng một vòng lặp và bước Newton trùng với nghiệm
đóng \eqref{eq:closed-form}. Đẳng thức \eqref{eq:newton-one-step} là hệ quả trực
tiếp của
\eqref{eq:derivatives}: Hessian không phụ thuộc $w$ nên xấp xỉ bậc hai mà Newton
dựng tại $w_0$ không phải xấp xỉ, nó chính là hàm mục tiêu.

Phép so có ý nghĩa ở đây do đó không phải số vòng lặp mà là chi phí một vòng.
Newton dùng toàn bộ Hessian nên nó phải lập ma trận Gram với chi phí $\bigO(nd^2)$
rồi phân rã với chi phí $\bigO(d^3)$, đặt cạnh $\bigO(nd)$ của một vòng gradient.
Với $n = \num{160000}$ và $d = 280$, số hạng $nd^2$ chi phối, và đó là lý do
Newton mất \SI{0.133}{\second} chứ không phải vài mili giây. Trên bài toán phi
tuyến, nơi Hessian thay đổi theo $w$ và mỗi bước Newton chỉ là một xấp xỉ địa
phương, số vòng lặp trở lại thành đại lượng đáng so.

\section{L-BFGS}
\label{sec:impl-lbfgs}

\begin{algorithm}[tp]
  \caption{Two-loop recursion for the L-BFGS direction}
  \label{alg:lbfgs}
  \begin{algorithmic}[1]
    \Require gradient $g_k$, memory of $m$ pairs $(s_j, y_j)$ with
             $\rho_j = 1 / y_j\trans s_j$
    \State $q \gets g_k$
    \For{$j = k-1, \dots, k-m$}
      \State $a_j \gets \rho_j \, s_j\trans q$; \quad $q \gets q - a_j y_j$
    \EndFor
    \State $r \gets \gamma_k \, q$
      \Comment{$\gamma_k = s_{k-1}\trans y_{k-1} / y_{k-1}\trans y_{k-1}$}
    \For{$j = k-m, \dots, k-1$}
      \State $b \gets \rho_j \, y_j\trans r$; \quad $r \gets r + (a_j - b) s_j$
    \EndFor
    \State \Return $-r$
  \end{algorithmic}
\end{algorithm}

L-BFGS không lập Hessian mà dựng xấp xỉ nghịch đảo của nó từ $m$ cặp hiệu gần
nhất $s_j = w_{j+1} - w_j$ và $y_j = \nabla f(w_{j+1}) - \nabla f(w_j)$, qua đệ
quy hai vòng ở thuật toán~\ref{alg:lbfgs}~\cite{nocedal2006numerical}. Mỗi vòng
của đệ quy chỉ gồm tích vô hướng và phép cộng vector nên toàn bộ hướng đi tính
được với chi phí $\bigO(md)$, cộng thêm $\bigO(nd)$ cho một lần lấy gradient.

Đó là lý do L-BFGS nằm giữa hai nhóm về cả chi phí lẫn tốc độ. Nó dùng nhiều
thông tin hơn gradient descent nên cần ít vòng lặp hơn, 8 thay vì 20 ở
bảng~\ref{tab:summary}, nhưng nó không phải trả cái giá $\bigO(nd^2)$ của Newton
vì nó không bao giờ chạm tới ma trận $d \times d$. Với $d$ lớn tới mức không lập
nổi ma trận ấy, đây là phương pháp duy nhất trong báo cáo còn chạy được.

Cả năm thuật toán đều đã qua bốn phép kiểm tra tính đúng đắn trước khi chạy lưới
tham số, mô tả ở phụ lục~\ref{app:tests}. Ngoài năm thuật toán này nhóm còn cài
Heavy ball, Newton-CG và Adam; phụ lục~\ref{app:cut} nói vì sao chúng không vào
phần so sánh chính.
